\subsection{Deng et al.\ (2023): CHGNet charge-informed universal MLIP}
\label{sec:appendix-example-deng2023}

\paragraph{Q1 -- Bibliographic identification.}
Deng, Zhong, Jun, Riebesell, Han, Bartel, and Ceder introduce CHGNet,
a charge-informed universal interatomic potential trained on the
Materials Project Trajectory Dataset (MPtrj). Published in
\emph{Nat.\ Mach.\ Intell.} 5, 1031--1041 (2023); arXiv:2302.14231.
\lstinputlisting[language=turtle]{artifacts/kg/listings/deng2023/q01-bibliographic.ttl}

\paragraph{Q2 -- Material system.}
The MLIP covers the September 2022 Materials Project release: 145{,}923
inorganic compounds spanning 94 elements. Modelled as a single
$\MaterialSystemC$ at the database-coverage level.
\lstinputlisting[language=turtle]{artifacts/kg/listings/deng2023/q02-material.ttl}

\paragraph{Q3 -- Reference calculation method.}
Reference data come from DFT calculations.
\lstinputlisting[language=turtle]{artifacts/kg/listings/deng2023/q03-reference-method-type.ttl}

\paragraph{Q4 -- Reference settings.}
DFT calculations use VASP with PAW pseudopotentials, PBE GGA (with
GGA+U on transition-metal atoms, e.g.\ $U=3.9$\,eV for Mn), a 520\,eV
plane-wave cutoff, and a reciprocal-space discretisation of 25
k-points/\AA$^{-1}$. The MP GGA/GGA+U mixing compatibility correction
is applied.
\lstinputlisting[language=turtle]{artifacts/kg/listings/deng2023/q04-reference-settings.ttl}

\paragraph{Q5 -- Training dataset.}
MPtrj contains 1{,}580{,}395 atomic configurations from 145{,}923 MP
compounds, with energies, forces, and stresses (plus 7{,}944{,}833 site
magnetic moments used as auxiliary supervision). Provenance is
\emph{published}.
\lstinputlisting[language=turtle]{artifacts/kg/listings/deng2023/q05-dataset.ttl}

\paragraph{Q6 -- Sampling strategies.}
Three strategies were combined: MP relaxation-trajectory sampling
(with extensive filtering for SCF convergence, energy bounds, and
StructureMatcher deduplication, then 8:1:1 train/val/test split by
\texttt{mp-id}); periodic-table chemical-space coverage (94 elements;
60+ elements with $>$100k site occurrences; 76 elements with
magnetic-moment annotations); and site-magnetic-moment decoration
(magmom auxiliary supervision acting as a charge-state regulariser
on the latent space).
\lstinputlisting[language=turtle]{artifacts/kg/listings/deng2023/q06-sampling.ttl}

\paragraph{Q7 -- MLIP method, components, implementation.}
CHGNet is a two-graph (atom + bond) message-passing network with
explicit angle features, atom/bond/angle convolution blocks, and
magmom regularisation of the latent atom features. Outputs are
energy, forces, stresses (auto-differentiated), and magnetic moments.
Loss is Huber ($\delta=0.1$) on all four with weights
$w_E=1$, $w_f=1$, $w_\sigma=0.1$, $w_m=0.1$. Optimisation is Adam with
CosineAnnealingLR. Implementation: CHGNet (PyTorch 1.12.0).
\lstinputlisting[language=turtle]{artifacts/kg/listings/deng2023/q07-method.ttl}

\paragraph{Q8 -- Hyperparameter settings.}
Reported settings: atom-graph cutoff $r_c = 5.0$\,\AA, bond-graph
cutoff $3.0$\,\AA, feature dimension 64, 4 convolution layers, batch
size 40, initial learning rate $10^{-3}$.
\lstinputlisting[language=turtle]{artifacts/kg/listings/deng2023/q08-hyperparameter-settings.ttl}

\paragraph{Q9 -- MLIP run and trained model.}
A single training run on MPtrj produces the pretrained universal
CHGNet (with-magmom variant, 400{,}438 trainable parameters).
\lstinputlisting[language=turtle]{artifacts/kg/listings/deng2023/q09-run-and-model.ttl}

\paragraph{Q10 -- Benchmark results.}
Test-set MAEs on the held-out MPtrj split (157{,}955 structures from
14{,}572 materials): 30\,meV/atom on energy, 77\,meV/\AA{} on force
components, 0.348\,GPa on stress components.
\lstinputlisting[language=turtle]{artifacts/kg/listings/deng2023/q10-benchmark-results.ttl}

\paragraph{Q11 -- Computational resources.}
\emph{Not reported.}
\lstinputlisting[language=turtle]{artifacts/kg/listings/deng2023/q11-resources.ttl}

\paragraph{Q12 -- Gaps.}
Beyond Q11, the most prominent gap is the \emph{charge-informed}
aspect of CHGNet: site magnetic moments are auxiliary outputs that
the ontology cannot represent as a first-class
$\dlConcept{CoveredProperty}$ (the model encodes them at a hyperparameter-
loss-weight level but not as a queryable target property). The
universal-coverage simplification of Q2/Q5 carries the same
limitation as for chen2022: 94 elements and many MP-task variants
collapse to a single $\MaterialSystemC$. Downstream applications
(LiMnO$_2$ MD, NaMnO$_2$ phase prediction, Li-ion conductor relaxations)
are also not encoded as separate $\BenchmarkResultC$ instances.
